\chapter{Serviços de comunicação e como
  solicitá-los}

\section{Solicitantes}

Os serviços de comunicação e divulgação de ações, temas técnicos,
científicos e políticos nas mídias institucionais do Conselho Regional
de Psicologia de São Paulo (CRP~SP) podem ser solicitados exclusivamente
pelo Plenário, pelas comissões permanentes e especiais e pelas Unidades
Administrativas (coordenações e assessorias), em conformidade com as
diretrizes desta Política de Comunicação.

O artigo 2º, inciso I, parágrafo 1º~da
\href{https://transparencia.cfp.org.br/wp-content/uploads/sites/29/2024/06/Resolucao-03-2024.pdf}{Resolução
  CRP nº~03/2024} determina que as ações propostas por subcomissões devem
ser aprovadas junto às comissões permanentes a que se encontrem
subordinadas. \textbf{A solicitação de serviços e produtos para as
  subcomissões, portanto, deve ser feita pelas comissões permanentes
  responsáveis.}

\subsection{Como solicitar serviços de
  comunicação}

Para requisitar serviços de comunicação, a solicitante deve:

\begin{enumerate}

  \item
        preencher e encaminhar à Unidade de Comunicação o formulário
        específico para a demanda, juntamente com toda a documentação prevista
        nos respectivos formulários. É essencial atender a todos os requisitos
        indicados para que a proposta seja comprovada;
  \item
        formalizar a solicitação de divulgação e/ou execução da atividade por
        meio de um processo no SEI!, dirigido à Unidade de Comunicação e com
        notificação por \emph{e-mail}. Esse setor irá prosseguir à tramitação
        e realizará as verificações cabíveis para viabilizar a solicitação
        e/ou a divulgação das atividades técnicas, científicas e políticas nos
        canais institucionais do CRP~SP.
\end{enumerate}

\section{Serviços}

\textbf{Atenção:} peças audiovisuais (\textbf{fotografias e gravações de
  voz, vídeo ou ambos}) destinadas à publicação nos canais oficiais do
CRP~SP devem \textbf{obrigatoriamente} ser acompanhadas de termo de
autorização para o uso de imagem devidamente preenchido por todas as
pessoas participantes das peças, além das seguintes informações (para
cada uma das peças):

\begin{enumerate}

  \item
        data e local de criação da imagem;
  \item
        descrição geral do assunto (nome de evento, campanha em que se insere
        etc.);
  \item
        nomes das pessoas gravadas, filmadas ou fotografadas.
\end{enumerate}

\subsection{Campanhas}

\textbf{Campanhas} são ações coordenadas que visam um objetivo comum:
por exemplo, conscientizar a categoria acerca de diretrizes éticas
emanadas do Sistema Conselhos.

As campanhas do CRP~SP têm por objetivo auxiliar na efetivação do
planejamento estratégico da autarquia. São solicitadas por comissões
permanentes e especiais e unidades administrativas.

Para toda campanha, deve ser elaborado, em parceria com a ComCom, plano
de comunicação que defina o escopo completo da campanha (mote, objetivo,
ações, linguagens, produtos, grade de veiculação etc.).

\subsubsection{Relações públicas e assessoria de
  imprensa}

A Unidade de Comunicação do CRP~SP media a comunicação
interinstitucional e atua junto à imprensa, sugerindo pautas e atendendo
a solicitações de informação e entrevistas, conforme a demanda interna e
externa.

Em toda a comunicação com outras instituições e, principalmente, com a
imprensa, a Unidade de Comunicação deve definir os meios e a abordagem a
serem empregados em parceria com a ComCom.

\subsubsection{Produção e apoio a
  eventos}

A Coordenação de Comunicação do CRP~SP é responsável pela gestão
(produção, cobertura e transmissão) de eventos, como mostras,
seminários, rodas de conversa, congressos, cinedebates, premiações e
eventos internos, promovidos pela autarquia \textbf{em nível estadual}.

A produção de eventos nos territórios é atribuição das respectivas
subsedes e deve ser solicitada por meio de
\href{https://forms.gle/yFMCSwpP6DQ5p9dA6}{formulário próprio}, cabendo
à Coordenação de Comunicação a orientação e a divulgação, além da
disponibilização das seguintes ferramentas digitais:

\begin{enumerate}

  \item
        inclusão do evento na Agenda CRP~SP;
  \item
        fornecimento de \emph{link} para inscrição;
  \item
        emissão de certificados.
\end{enumerate}

\paragraph{Cobertura jornalística de
  eventos}

A unidade de comunicação deve providenciar a cobertura, na forma de
fotos e vídeos, para eventos de caráter estadual. Dos eventos em
territórios, a subsede de referência é responsável pela cobertura.

Em qualquer caso, a divulgação de balanços (cobertura pós-evento) fica
condicionada ao preenchimento de
\href{https://forms.gle/aZUAWsRiPbCLaaUL7}{relatório próprio} por
profissional qualificado para abordar o assunto de que tratou o evento e
à disponibilidade de fotografias, que devem trazer, no próprio arquivo
da imagem ou em arquivo de texto complementar:

\begin{enumerate}

  \item
        data de criação da imagem;
  \item
        descrição geral do assunto (nome de evento, campanha em que se insere
        etc.);
  \item
        nomes das pessoas gravadas, filmadas ou fotografadas.
\end{enumerate}

\subsection{Serviços editoriais}

\subsubsection{Revisão de textos}

Os serviços editoriais prestados pela Unidade de Comunicação do CRP~SP
são:

\begin{enumerate}

  \item
        revisão de texto;
  \item
        preparação de originais;
  \item
        criação de identidade visual e projeto gráfico;
  \item
        execução de projeto gráfico;
  \item
        produção gráfica/produção \emph{web}.
\end{enumerate}

O serviço de \textbf{revisão} compreende a correção de textos de acordo
com as normas gramaticais e ortográficas vigentes, além de adequação à
linguagem gendrada e demais padrões editoriais do CRP~SP. Pode, ainda,
envolver alterações no léxico, na sintaxe e na estrutura do texto, com o
objetivo de torná-lo mais acessível e/ou mais alinhado à \emph{persona}
do Conselho.

A \textbf{preparação de originais} envolve, além do serviço de revisão,
a organização de documentos mais extensos visando sua publicação na
forma de monografia (livro, apostila ou similar).

Para que os documentos produzidos possam ser publicados, é necessária a
criação de uma \textbf{identidade visual} que reflita os objetivos da
publicação e de um \textbf{projeto gráfico} subordinado a esta
identidade visual que garanta um bom padrão estético e de legibilidade.
A execução desse projeto gráfico é o que costumamos chamar de
diagramação.

A etapa seguinte à diagramação é a \textbf{produção gráfica}, e vai
desde a escolha do material e da forma de impressão à contratação de
fornecedores e à avaliação do produto final.

\subsection{Formulários e suas
  funções}

\subsubsection{Formulário de solicitação e regramento de divulgação nas
  mídias do
  CRP~SP}

Este formulário destina-se às comissões permanentes e especiais e às
unidades administrativas que desejam solicitar divulgação institucional
nas publicações oficiais do Conselho Regional de Psicologia de São Paulo
(CRP~SP). Todas as propostas serão submetidas a uma avaliação prévia
pela Unidade de Comunicação e pela Comissão de Comunicação Institucional
(ComCom), garantindo o alinhamento com as diretrizes institucionais e a
coerência com os objetivos do CRP~SP.

O formulário deve ser utilizado para solicitar a divulgação de:

\begin{enumerate}

  \item
        atividades técnicas, científicas e políticas;
  \item
        divulgação de temas abordados pelo Centro de Referências Técnicas em
        Psicologia e Políticas Públicas (Crepop);
  \item
        divulgação de ações de entidades parceiras;
  \item
        nota de pesar;
  \item
        nota de posicionamento;
  \item
        informes administrativos e financeiros;
  \item
        orientações.
\end{enumerate}

Esses conteúdos serão direcionados aos canais institucionais com o
objetivo de informar a categoria sobre representações e ações do
Conselho.

\subsubsection{Formulário de solicitação de
  evento}

\textbf{Link:
  \href{https://forms.gle/yFMCSwpP6DQ5p9dA6}{https://forms.gle/yFMCSwpP6DQ5p9dA6}}

O conteúdo deste formulário deve ser anexado em formato PDF ao processo
aberto no SEI, juntamente com o ``Formulário de aprovação de despesas
com atividades precípuas''. Após a aprovação no SEI, envie o número do
processo, o formulário de evento atualizado e a descrição das demandas
para a Unidade de Comunicação, setor de eventos.

\subsubsection{Relatório final do
  evento}

\textbf{Link:
  \href{https://forms.gle/aZUAWsRiPbCLaaUL7}{https://forms.gle/aZUAWsRiPbCLaaUL7}}

O Relatório Final do Evento do CRP~SP é um documento que reúne
informações essenciais sobre eventos, ações e atividades orientativas
promovidas pelo Conselho. Este relatório inclui detalhes como
identificação do evento, objetivos, atividades realizadas, avaliação dos
resultados e registros multimídia. Seu objetivo é documentar de forma
organizada e eficiente a natureza e o impacto do evento, conveniente
para fornecimento de contas, avaliações futuras e possível divulgação
nos canais do CRP~SP.

Esses procedimentos visam garantir a consistência e a efetividade das
ações de comunicação, além de garantir o alinhamento com a missão e os
valores do CRP~SP.